\begin{algorithm}[!ht]
\caption{Light TSMixup: background generation from a sinusoidal frequency pool}
\label{alg:light_tsmixup}
\begin{algorithmic}[1]
\Require Frequency pool $\mathcal{P}=\{f_1,\ldots,f_M\}$ in Hz, with $M=574$, built from the design as described in this appendix; maximum components $K=10$; symmetric Dirichlet concentration $\alpha=1.5$; draw length $l$, the probing context plus the forecast horizon; sampling frequency $f_s$.
\Ensure One background realisation $x_{1:l}$ of unit variance.
\For{$m \gets 1$ to $M$}
    \State $x^{(m)}_{1:l} \gets \sin\!\left(2\pi f_m\, n / f_s\right),\ n=0,\ldots,l-1$ \Comment{one sinusoid per pool frequency}
\EndFor
\State $k \sim U\{1,K\}$ \Comment{number of components to mix}
\For{$i \gets 1$ to $k$}
    \State $m \sim U\{1,M\}$ \Comment{draw a source, with replacement}
    \State $\tilde{x}^{(i)}_{1:l} \gets x^{(m)}_{1:l} \Big/ \left(\frac{1}{l}\sum_{j=1}^{l}\left|x^{(m)}_j\right|\right)$ \Comment{apply mean scaling}
\EndFor
\State $[\lambda_1,\ldots,\lambda_k] \sim \operatorname{Dir}([\alpha_1=\alpha,\ldots,\alpha_k=\alpha])$ \Comment{sample mixing weights}
\State $x_{1:l} \gets \sum_{i=1}^{k}\lambda_i\tilde{x}^{(i)}_{1:l}$ \Comment{take a weighted combination of scaled sources}
\State \Return $x_{1:l} / \operatorname{std}\!\left(x_{1:l}\right)$ \Comment{unit variance, so the injected tone sets a true SNR}
\end{algorithmic}
\end{algorithm}


\medskip 

Light TSMixup is a variant of the TSMixup augmentation described in the Chronos data-generation procedure\autocite{ansariChronosLearningLanguage2024}. The original algorithm samples subsequences from a collection of real training datasets, applies mean scaling to each sampled series, draws mixing coefficients from a symmetric Dirichlet distribution and returns their convex combination. The variant used here preserves that stochastic mixing structure and replaces the real dataset collection with a pool of sinusoidal sources of known frequency, so that every generated series remains a superposition of known components and no loss measured at a predicted site can be attributed to unknown spectral content in the background.

\subsection{The source pool}

The pool holds $397$ unit-amplitude sinusoids, one per frequency, and is built from the design rather than fixed by hand: adding a geometry to \autoref{tab:hfModels} adds its own frequencies automatically. It spans $2$ to $150$\,Hz, which is not the analysed band but a deliberate restriction of it, discussed at the end of this section. It has four parts.

The first is every integer frequency from $2$ to $150$\,Hz, which is $149$ sources. Covering the integer grid uniformly is what keeps the pool neutral with respect to the hypotheses: a candidate frequency is represented in the background exactly as the frequencies flanking it are, so the background cannot by itself make a candidate easier or harder to recover than its controls.

The second is the nine members of $\mathcal F_{\mathrm{lock}}$ below $150$\,Hz that are not integers, which the first part cannot reach. They arise in the geometries whose stride or patch size is $12$, $20$ or $24$, where $f_s/D$ is not a whole number of hertz: the stride comb of $S=12$ falls at $42.\overline{6}$ and $85.\overline{3}$\,Hz, and those of $S=20$ and $S=24$ at multiples of $25.6$ and $21.\overline{3}$\,Hz. Without them the pool would represent the integer candidates and omit the rest.

The third is the control frequencies $f_k\pm\delta$ of every candidate that falls in the pool's range, $141$ distinct values of which $89$ are not already integers. They matter for the same reason the integer grid does, and more sharply. The offset $\delta$ is chosen per site as the largest one that keeps both controls clear of $\mathcal F_{\mathrm{lock}}$, so it is rarely a whole number of hertz: without this group every locked arm of a contrast triplet would sit on a pool frequency while fewer than half the control arms would, and the background would carry a little more energy on one side of \autoref{eq:d2contrast} than on the other. Since that contrast assumes its three arms differ only in frequency, the asymmetry would be indistinguishable from the effect $H1$ is meant to detect. With the controls in the pool both arms are covered alike.

The fourth is $150$ further non-integer frequencies spread evenly across the same range, each held at least $0.2$\,Hz away from every candidate and from the integer grid. They give the background content that no geometry predicts, which is what prevents the pool from being a description of $\mathcal F_{\mathrm{lock}}$ and its controls and nothing else.

\subsection{Generation}

One realisation is produced by drawing a number of components $k$ uniformly between one and the maximum $K$, drawing that many sources from the pool with replacement, scaling each by the reciprocal of its mean absolute value, and returning their convex combination under weights drawn from a symmetric Dirichlet distribution of concentration $\alpha$. The values adopted are $K=10$ and $\alpha=1.5$, and every source is generated at the full draw length, the probing context plus the forecast horizon, so no source is tiled or truncated. Drawing with replacement means a frequency may enter one mixture more than once, and a mixture therefore carries at most ten of the $397$ frequencies: a single background is sparse in that range even though the pool is dense in it.

Two consequences follow. The concentration $\alpha$ sets how evenly the mixing weights are spread, so values near zero return something close to a single sinusoid while larger values return balanced mixtures; at $\alpha=1.5$ the components are comparable in size without being equal. And because the sources are defined in hertz rather than in cycles per sample, the pool is tied to the sampling rate, hence to the lock arithmetic of \nameref{sec:two}: the same construction at a different $f_s$ yields a different pool, which is the intended behaviour.

The pool stops at $150$\,Hz while the analysed band runs to $250$\,Hz, and the asymmetry is deliberate. Above that limit the reconstruction of \nameref{sec:appSweep} returns almost nothing for any geometry, so a background component there would be spectral content the model cannot carry in either arm of a contrast, and the mixture would spend its ten components on frequencies that cannot inform $H1$. What the restriction costs is stated plainly: $87$ of the candidate frequencies of the design lie above $150$\,Hz and are probed on a background that is silent there. That silence is not a bias, because a candidate and its two controls lose the background together and the contrast of \autoref{eq:d2contrast} is paired, but it does mean the claim that the background covers the analysed band holds only up to $150$\,Hz.

The probe tone is not part of this generator. A realisation is normalised to unit variance and the tone is added afterwards, at the frequency and phase the design requires and at a fixed signal-to-noise ratio, so the background and the quantity under test are produced independently of one another.

\autoref{alg:light_tsmixup} states the procedure step by step, with the values adopted here; the pool it takes as input is the one described in this appendix.